% Auto-generated proof export skeleton for TEXTBOOK-PART-IV-CHAPTER-15-MINIMAX-LOWER-BOUNDS-SPINE
\section{Chapter 15 Lemma 15.1, Theorem 15.2, and Exercise 15.7 data-processing leaf}

\paragraph{Status.}
This note is synchronized from locally compiled declarations. Lemma~15.1 and
Theorem~15.2 are compiled; Exercise~15.7 remains partial.

\paragraph{Lean declarations.}
\texttt{LowerBounds.banditHistoryRelativeEntropy\_eq\_expectedPulls\_sum},
\texttt{LowerBounds.finiteArmedGaussianMinimaxLowerBound},
\texttt{LowerBounds.unitGaussianMinimaxExpectedPseudoRegret\_ge},
\texttt{LowerBounds.klDiv\_map\_le}, and
\texttt{LowerBounds.klDiv\_observedBanditHistory\_le\_expectedPulls\_sum}.

\paragraph{Proof map.}
Expanding the canonical history law one round at a time cancels the common
randomized policy kernel. The remaining conditional reward divergence is
averaged under the first environment and regrouped by arm, proving Lemma~15.1
with its directed KL and first-law expected pull counts.

For Theorem~15.2, the least-explored alternative and the unit-Gaussian KL
formula bound the base-to-changed history divergence by $1/2$ at the source
gap. Bretagnolle--Huber applied to the source pull-count event forces one of
the two expected pseudo-regrets above
$\sqrt{(k-1)n}/27$. The Lean endpoint constructs the unit-cube witness for
every randomized policy and proves the worst-case/minimax consequence.

The generic data-processing leaf identifies the mapped Radon--Nikodym density
with a conditional expectation and applies conditional Jensen to the convex
KL integrand, with a separate infinite-KL branch. Its deterministic-history
corollary follows from Lemma~15.1. It does not prove Exercise~15.7: the stopped-
history information bound and factorization of an arbitrary
$\mathcal F_\tau$-measurable random element remain planned.
